\documentclass[11pt,a4paper]{article}
\usepackage[T1]{fontenc}
\usepackage{lmodern,amsmath,amssymb,amsthm,mathtools,booktabs,microtype,xcolor}
\usepackage[margin=26mm]{geometry}
\usepackage[colorlinks=true,linkcolor=blue!45!black,citecolor=blue!45!black,urlcolor=blue!45!black]{hyperref}
\newtheorem{theorem}{Theorem}[section]\newtheorem{lemma}[theorem]{Lemma}
\newtheorem{proposition}[theorem]{Proposition}
\newcommand{\Q}{\mathbb Q}\newcommand{\Z}{\mathbb Z}\newcommand{\HH}{\mathcal H}
\DeclareMathOperator{\Tr}{Tr}\DeclareMathOperator{\Nm}{Nm}
\setlength{\emergencystretch}{2em}
\title{A family of defining quartics in 279 variables}
\author{}\date{9 November 2026}
\begin{document}\maketitle
\begin{abstract}
Under the stated arithmetic premises of the elliptic source, we specify
an integer quartic family in 279 total variables whose sufficiently
large members define the integers in the rationals. Four disjoint
translates of the contact set give a height exponent $19/4$. This permits
quadratic Taylor approximations to the formal logarithm, with comparison
orders twelve and three. Even valuations of reciprocal abscissas allow
sixteen local tests to be represented by five factors. The complete
quadratic and expanded coefficient tables are independently verified.
The family has 3692 monomials for every positive parameter, and exact
coefficient height $80989^{4N+2}+1$ for $N\ge23$. The sufficient defining
parameter remains ineffective because the aggregate-trace height theorem
has an ineffective additive constant.
\end{abstract}

\section{The family and its arithmetic scope}
All variable counts include the distinguished input $X$. For a fixed
integer $N\ge1$ let $A_N=80989^{2N+1}$. The complete tables
\nolinkurl{quadratic-family279.json} and
\nolinkurl{polynomial-family279.json} specify 253 primitive integer
quadratic polynomials and their sum of squares
\[
 P_N(X,Y_1,\ldots,Y_{278})=\sum_{j=1}^{253}Q_{N,j}^2.
\]
The parameter $A_N$ is a coefficient, not a variable.

\begin{theorem}\label{thm:main}
Assume the arithmetic premises stated below. There is a fixed integer
$N_0\ge27464$ such that every $N\ge N_0$ satisfies
\[
 x\in\Z\quad\Longleftrightarrow\quad
 \exists\mathbf y\in\Q^{278}\ P_N(x,\mathbf y)=0
 \qquad(x\in\Q).
\]
For every $N\ge1$ the polynomial $P_N$ has degree exactly four,
content one, and 3692 nonzero coordinate monomials. For $N\ge23$,
\[
                       H(P_N)=80989^{4N+2}+1.
\]
The construction specifies every member exactly. The proof does not
determine a sufficient numerical value of $N_0$.
\end{theorem}

Set
\[
 K=\Q(\rho),\quad \rho^2=2,\quad D=(3-\rho)^2=11-6\rho,
 \qquad E:y^2=x^3-Dx,
\]
\[
 P=2(1+\rho,2\rho)
 =\left(\frac{57-28\rho}{8},\frac{-660+417\rho}{32}\right).
\]
We retain the following arithmetic premises, with the precise domains
and hypotheses recorded in \cite{Dependencies,Source,Weights}:
\begin{enumerate}
\item $2E(K)=\Z P$ and $h(x([n]P))\ge n^2-6$ for every nonzero
integer $n$, together with the specified point charts and group domains.
\item The five-contact height estimate with exponent nineteen,
exceptional cutoff $76176$, and the fixed constant
\[
 H=2^{110}210^{1800}\left(\prod_{p\le76176}p\right)^{19}+128.
\]
\item The pole obstruction outside $2,3,17$, and its simultaneous
realization for positive inputs less than $17$ with square reduced
denominator and the stated local squareclasses and valuation bounds
at those three primes. Its normalized quadratic equations are retained.
\item The simultaneous prime-pattern, conjugate reduction-order,
finite congruence, and real approximation assertions used by the source.
The weighted selection and fixed denominator dilation are justified
from the cited analytic estimates in \cite{Selection}.
\end{enumerate}
The aggregate trace-height theorem used below is the proved consequence
of the family height theorem described in \cite{Aggregate}; its additive
constant is ineffective. The contact amplification in \cite{Contacts}
is an elementary consequence of the retained five-contact estimate,
not an additional arithmetic assumption. Exact coefficient verification
does not prove the retained arithmetic premises.

\section{Four contact blocks and the two scales}
Put
\[
 \Theta=\{0,-1/2,-1,1/2399,-32/407\},\qquad
 \Theta_4=\bigcup_{j=0}^3(\Theta+3j),\qquad
 F_4(t)=\prod_{\theta\in\Theta_4}(t-\theta).
\]
Let $\Sigma$ consist of the primes at most $76176$ and the two primes
$123121,161303$. Outside this set, all contact differences and all
numbers $1-\theta$ are units. The contact radical $M_4(t)$ is the
product of primes $p\notin\Sigma$ for which
$v_p(t-\theta)$ is positive and odd for some $\theta\in\Theta_4$.
The exact contact certificate in \cite{Contacts} verifies this
exceptional set, including the normalizing value $t=1$.

Let $E_4=123121\cdot161303=19859786663$. The contact amplification
 theorem gives
\begin{equation}\label{eq:contactheight}
 h(t)\le\log26+HE_4^{19}M_4(t)^{19/4}
       \le H_4M_4(t)^{19/4},\qquad
 H_4=\log26+HE_4^{19},
\end{equation}
for $t\notin\Theta_4$, and $M_4(1)=1$.
For clarity, its proof applies the five-contact estimate to the four
maps $2+1/(t-3j)$, each with determinant $-1$. Their contact preimages
are disjoint outside $\Sigma$. Each map changes logarithmic height
by at most $\log26$. Multiplication of the four resulting inequalities
and disjointness of the radicals gives \eqref{eq:contactheight}.
The full treatment of fixed exceptional primes is in \cite{Contacts}.

The retained root points are
\[
 Q_1=[n]P,\quad Q_2=[a]P,\quad Q_3=[t]P,\quad Q_4=[b]P,
 \qquad B=\max(|n|,|a|,|t|,|b|),
\]
with nonzero integer labels. Set
\[
 T=\frac12\Tr_{K/\Q}\left(\sum_{i=1}^4x(Q_i)\right).
\]
It is positive, and the aggregate trace theorem gives fixed constants
$M_1>0$, $M_2\ge1$ such that
\begin{equation}\label{eq:aggregate}
                         B^2\le M_1h(T)+M_2.
\end{equation}
Use the two nonzero rational square scales $a_*=1,T^2$.
For each scale define, on its nonzero domain,
\begin{equation}\label{eq:detector}
 t_*=a_*\lambda,\quad \delta=t_*-1,\quad
 R=\frac{F_4(t_*)}{a_*^{20}\delta^{22}},\quad
 e=\frac{(a_*\delta)^{20}}{F_4(t_*)},\quad
 q=\frac1{R\delta^4}.
\end{equation}
The actual quadratic system clears these expressions by
\begin{equation}\label{eq:clearing}
             F_4(t_*)e=(a_*\delta)^{20},\qquad q\delta^2=e.
\end{equation}
In particular $e\delta^2=R^{-1}$ when $\delta\ne0$.
No contact value can occur, since $F_4(t_*)=0$ would contradict the
first equation. If $q=0$, then $e=0$, and that equation forces
$\delta=0$. At $\delta=0$ the contact radical is one; no reciprocal
is used on this stratum.

If both radicals divide a nonzero integer $J$ to order $k$, then
$2h(T)\le h(\lambda)+h(T^2\lambda)$ and
\eqref{eq:contactheight}--\eqref{eq:aggregate} imply
\begin{equation}\label{eq:heighttransfer}
 B^2\le C|J|^{19/(4k)}+M_2,\qquad C=M_1H_4.
\end{equation}
The two radicals need not be coprime for this implication.

\section{Local tests and the approximation denominators}
Let $\sigma$ be conjugation, $f=1/x$, and $z=-x/y$, extended by zero
at the identity where these functions are regular. The retained
point relations are
\[
 U_0=H,\quad V_0=H-P,\quad U_i=[u_i]P,\quad V_i=U_i-Q_i.
\]
The charts permit $U_i=O$ or $V_i=O$ for $i>0$.
Put
\[
 G_2(w)=1-\frac25w-\frac2{15}w^2,
 \quad \ell_0=z(V_0)G_2(Df(V_0)^2),\quad
 L_i=\frac{z(U_i)}{z(H)},
\]
\[
 L_1'=\sigma\!\left(\frac{z(V_1)}{\ell_0}\right),\qquad
 L_i'=\sigma\!\left(\frac{z(V_i)G_2(Df(V_i)^2)}{\ell_0}\right)\quad(i>1).
\]
These quotients are interpreted only after the denominator exclusions
below. Write
\[
 D_H=f(H),\quad A_i=f(U_i),\quad B_i=\sigma f(V_i),\quad
 C_i=L_i-L_i',\quad I=L_1-\beta,\quad \beta=51X^2+2504,
\]
\[
 G=L_3^2-L_2(L_2-2)L_3L_4+(L_2-2)^2L_4^2-1.
\]
The required orders are six for $D_H$, twelve for $C_2,C_3,C_4,G$,
three for $C_1,I$, and one for the nine $A_i,B_i$.

\begin{lemma}\label{lem:g2}
The polynomial $G_2$ has no zero in $K$.
\end{lemma}
\begin{proof}
Its roots are $(-3\pm\sqrt{39})/2$, and $39$ is neither a rational
square nor twice a rational square. Explicitly, if
$(u+v\sqrt2)^2=39$ with $u,v\in\Q$, then $uv=0$; the two cases
would require $u^2=39$ or $v^2=39/2$, both impossible.
The independent finite certificate also verifies that the integer
numerator $15-6w-2w^2$ has no root modulo 17, a prime that splits in
$K$. Its leading coefficient is a unit there, so a hypothetical local
root would be integral and give a forbidden residue root.
\end{proof}

The finite $H$ chart is $xf=1$, $S^2=x-Df$, $H=(x,xS)$.
The source group calculation writes $H=[m]R_0+T_0$ with $P=2R_0$,
$m\in\Z$, and $T_0\in E[2](K)$. Put $h=m/2$ for its rational
free index. The exceptional classes $E[2](K),\pm P+E[2](K)$ are
excluded before any logarithmic ratio is interpreted. Indeed their
finite allowed reciprocal abscissas are units outside the fixed small
primes, whereas extraction below would require positive valuation
at any contact. All radicals would then be one, forcing
$B^2\le C+M_2$ by \eqref{eq:heighttransfer}. The retained scalar
lower bound in Section~\ref{sec:soundness} contradicts this for every
sufficiently large parameter. The explicit torsion translations and
norm values supporting the unit assertion are those recorded in
\cite[Section 3]{Even}. This exclusion uses no comparison quotient.

Consequently $m\ne0,2$, and $H,V_0$ are nontorsion. Their required
nonzero abscissas and ordinates, and all subtraction denominators,
follow from the retained chart identities. In particular
$z(H),z(V_0)\ne0$. Lemma~\ref{lem:g2} gives $\ell_0\ne0$ globally.
The analytic formal logarithms have a separate nonvanishing proof
from the formal logarithm isometry when the points are formal.

\section{Five factors and extraction at one place}
At a contact prime $p\notin\Sigma$, $\delta$ is a unit and
$\nu=v_p(R)$ is odd; it need not initially be positive. All valuations
of $K$ are unramified and normalized by $v(p)=1$.
For a point on $E$, every finite value of $v(f)$ is even: if
$v(x)<0$ then $2v(y)=3v(x)$, while if $v(x)>0$ then
$2v(y)=v(x)$; the remaining case is valuation zero.
This also holds for conjugated reciprocal abscissas.

Use
\begin{align*}
 T_1&=1+q^{23}C_2^2+q^{10}D_H^2+q^4A_3^4+q^8A_4^8,\\
 T_2&=1+q^{23}C_3^2+q^{10}C_1^4+q^4B_0^4,\\
 T_3&=1+q^{23}C_4^2+q^2A_1^2+q^4B_2^4,\\
 T_4&=1+q^{23}G^2+q^2A_2^2+q^4B_3^4,\\
 T_5&=1+q^5I^2+q^2B_1^2+q^4B_4^4,
 \qquad W=T_1T_2T_3T_4T_5.
\end{align*}

\begin{proposition}\label{prop:extract}
At each contact prime, the retained pole system and the new cleared
trace equations imply $\nu>0$ and give all sixteen required test
orders at one place above that prime.
\end{proposition}
\begin{proof}
In every factor the finite nonconstant valuations in successive
positions have exact powers of two divisors $1,2,4,8$, omitting the
last when absent. For example the first factor has valuations
\[
 2v(C_2)-23\nu,\quad 2(v(D_H)-5\nu),\quad
 4(v(A_3)-\nu),\quad 8(v(A_4)-\nu).
\]
Each parenthesis that is required to be odd is odd, by evenness of
the reciprocal-abscissa valuations. The second position involving
$C_1$ is $2(2v(C_1)-5\nu)$ and is likewise twice an odd integer.
The first position of $T_5$ is $2v(I)-5\nu$, which is odd.
Thus the minimum, including the constant term of valuation zero,
is attained once. Every factor is nonzero and has valuation at most
zero, and a factor of valuation zero is one in the residue field.

Write $W=w_0+w_1\rho$. The literal cleared trace equations are
\[
 w_0^2-2w_0\tau-2w_1^2=0,\qquad e\delta^2\mathfrak b=\tau^2.
\]
Since $W\ne0$, the first equation gives $w_0\ne0$ and
$\tau=1/\Tr(W^{-1})\ne0$. Therefore
$\mathfrak b=R/\Tr(W^{-1})^2$. If
$\ell=\max_{v\mid p}v(W)\le0$, the trace inequality gives
$v_p(\mathfrak b)\le\nu+2\ell$.
This is an odd valuation; the pole obstruction forces it positive.
Hence $\nu>0$ and $\ell>-\nu/2$.
At a place attaining $\ell$, every term of every factor has valuation
strictly greater than $-\nu/2$, because the factor valuations are
nonpositive and add to $\ell$.

The first-column comparisons satisfy
$v(C)>(23/2-1/4)\nu\ge11.25$, hence order twelve.
For $I$, the bound is $v(I)>(5/2-1/4)\nu\ge2.25$, hence order three.
The $C_1$ term gives $v(C_1)>(10/4-1/8)\nu\ge2.375$, also order
three. For $D_H$ it gives $v(D_H)>(5-1/4)\nu\ge4.75$;
evenness makes this at least six. Every order-one term gives
$v(f)>(1-2^{-j-1})\nu>0$ in its position $j\ge1$.
Zero tests automatically satisfy the asserted inequalities.
\end{proof}

\subsection{Integrality precedes the depth estimate}
The extracted positive $f$ valuations put $H,U_i$ on the primary
formal branch and $V_0,V_i$ on the conjugate formal branch.
At a nonidentity formal point of depth $r=v(z)\ge1$, write
\[
 \log_E(Q)=z(Q)\mathcal A(Q),\qquad
 \mathcal A(Q)=G(Df(Q)^2),\qquad \mathcal A(O)=1.
\]
The formal series satisfies
\[
 G(w)=1-\frac25w-\frac2{15}w^2+O(w^3),\qquad
 g_j=\frac{4j-5}{4j+1}g_{j-1}.
\]
The coefficient bound and full-tail proof in \cite[Lemma 3.1]{Even}
give $\mathcal A\equiv1\pmod{p^{4r}}$,
$G_2/\mathcal A\equiv1\pmod{p^{12r}}$, and
$1/\mathcal A\equiv1\pmod{p^{4r}}$.
These analytic units equal one at $O$, so zero numerator points
require no quotient $\log_E(O)/z(O)$.

The integral group identities
\[
 [m]U_i=[2u_i]H,\qquad
 [m-2]V_i=[2(u_i-n_i)]V_0
\]
kill the torsion component of $H$ and yield full logarithm ratios
$r_i=u_i/h$ and $r_i'=(u_i-n_i)/(h-1)$.
At the odd prime, $h,h-1$ are integral and at least one is a unit.
Hence at least one full ratio is integral. Multiplication by the
analytic units preserves the corresponding approximate valuation.
The positive-order cross test transfers integrality to the other
approximate ratio and thus to both full ratios.

Only now use $L_i=z(U_i)/z(H)$ integral. The $H$ test gives
$v(z(H))\ge3$, so every nonzero $U_i$ has depth at least three.
The primary errors have order at least twelve. The alternate errors
have order at least four for $i=1$ and twelve for $i>1$.
The integral full ratios convert these relative bounds into absolute
congruences. Multiplication by $h,h-1$ and subtraction then give
\[
 L_1\equiv n\pmod{p^3},\qquad
 L_i\equiv n_i\pmod{p^{12}}\quad(i>1).
\]
This argument includes zero $U_i$ or $V_i$, half-integral $h$, and
inert primes. At the inert primes conjugation preserves the valuation.

\section{Global soundness for sufficiently large parameters}\label{sec:soundness}
The retained scalar circuit gives, for $c=|a-1|$,
\begin{equation}\label{eq:scalar}
 c\ge80989^{L_N},\qquad
 L_N=\left\lfloor\frac{187(2N-1)}8\right\rfloor+1,
 \qquad h(\beta)\le\HH c^{19/24}.
\end{equation}
These statements, including their integral group identities and
exceptional-point exclusions, are the scalar conclusions of
\cite{Source}; no comparison logarithm is used to establish them.
Their circuit is retained, including the actual positive contact
product $r$ and the normalizer
\[
 r=\beta(\beta-1)(\beta-2401)(32\beta+343)(2\beta-1),
 \qquad R_s=\frac{A_Nr}{1+80989^2r}.
\]

Let
$\Delta=t^2-a(a-2)tb+(a-2)^2b^2-1$.
The local comparison shows that each twelfth power of a comparison
radical divides $\Delta$. If $\Delta\ne0$, then
$|\Delta|\le14B^4$, and \eqref{eq:heighttransfer} gives
\[
                  B^2\le C14^{19/48}B^{19/12}+M_2.
\]
Since $19/12<2$, this bounds $B$ by a fixed constant. Equation
\eqref{eq:scalar} gives $B\ge c-1\to\infty$ uniformly as
$N\to\infty$, a contradiction. Thus $\Delta=0$ eventually.
The retained integral Lucas descent then gives
\[
                         \log B>|a|-3\ge c-4.
\]

If $n-\beta\ne0$, let $J$ be its nonzero reduced numerator.
The local input comparisons make both cubed radicals divide $J$,
while $|J|\le2B\exp(h(\beta))$. For $B^2\ge2M_2$,
\eqref{eq:heighttransfer} gives
\[
 \log B\le\frac{12}{5}\log(2C)+\frac{19}{5}\log2
                +\frac{19}{5}h(\beta)
 \le C'+\frac{19}{5}\HH c^{19/24}.
\]
The right side is $o(c)$, contradicting the preceding lower bound.
Hence $n=\beta\in\Z$. If $X=u/v$ is reduced with $v>0$, then
$v^2\mid51$, and squarefreeness of $51$ gives $v=1$.
All bounds depend only on fixed constants. Since $L_N$ is increasing,
a single $N_0$ works for every larger $N$. Its dependence on the
aggregate additive constant is ineffective.

\section{Simultaneous completeness}
Fix an integer input and the coefficient parameter. Retain the scalar
and Lucas construction of the source, and choose the integer labels,
root points, two scales, and compact real and finite-place neighborhoods
before selecting primes. The neighborhoods avoid every retained
rational denominator; Lemma~\ref{lem:g2} ensures that the new
denominator introduces none. Every local test then has a fixed lower
valuation bound at each prime in the finite preparation set.

Choose a fixed positive integer $D_0$ clearing all contact coefficients
$D_0^2\theta/a_*$ and sufficiently divisible at the prepared primes,
including $\Sigma$. In the prime pattern write
\[
                         \lambda=\frac{u}{D_0^2w^2}.
\]
The zero contact makes $u$ one of the selected primes, outside that
finite set. Thus $v_p(\lambda)\le-2v_p(D_0)$ there. Make this upper
bound low enough that $t_*=a_*\lambda$ has valuation below zero and
all nonzero contacts. The exact valuation formulas are then
\[
 v_p(F_4(t_*))=20v_p(t_*),\quad v_p(e)=20v_p(a_*),\quad
 v_p(q)=18v_p(a_*)-2v_p(\lambda).
\]
The last expression becomes uniformly arbitrarily large. All new
weighted test terms are therefore small simultaneously at the fixed
prepared places. This is a fixed denominator dilation of the same
finite prime pattern, allowed by \cite{Selection}.

Select the enlarged finite set of contact primes with the inherited
opposite-order separation. The common-index congruences may retain
primary depth twenty for $H$ and the $U_i$; they are stronger than
needed here. At the opposite places the points are formal. The new
Taylor approximations give comparison precision twelve, and the
constant approximation at $V_1$ gives precision four. Thus every
nonconstant term of every factor has positive valuation at the chosen
primary place, where $v_p(R)=1$. There $W\equiv1$. At the opposite
place, the valuation-class argument gives either positive valuation
for $W^{-1}$ or residue one. Its trace is consequently a unit with
residue one or two, including the inert case. Hence the pole input
has valuation one at every selected contact prime.

The exact prime-pattern identity is
\[
 R=\frac{\displaystyle\prod_{\theta\in\Theta_4}
          (u-(\theta/a_*)D_0^2w^2)}
        {(D_0^2w^2)^{20}\delta^{22}}.
\]
Its denominator is a square, so at every prime outside the selected
prime numerators its valuation is even. Division by the squared trace
preserves this parity. This argument includes primes dividing $\delta$;
no unit assertion about $\delta$ is made away from contacts.
The chosen congruences give the local squareclasses at $2,3,17$,
and the preparatory divisibility gives the required valuation bounds.

On the fixed compact real neighborhoods,
\[
 R\sim a_*^{-22}\lambda^{-2},\qquad
 q\sim a_*^{18}\lambda^{-2},\qquad
 W\longrightarrow1,\qquad\Tr(W^{-1})\longrightarrow2.
\]
All test powers in the factors are even, so the pole input tends to
zero through positive values. It eventually lies in $(0,17)$.
The simultaneous selection supplies the required real size after all
fixed data have been chosen. The pole-realization input supplies its
witnesses, and the exact power identities fill the quadratic circuit.
This proves completeness with no additional condition on $N$.

\section{Exact quadratic specification}
The weights are given by seven scalar quadratic rows
\[
 q_2=q^2,\quad q_4=q_2^2,\quad c q_4=q,\quad q_8=q_4^2,
 \quad q_5=qq_4,\quad q_{10}=q_5^2,\quad c q_{23}=q_{10}^2.
\]
At every contact $q\ne0$, so these are the indicated powers, including
$c=q^{-3}$. Their zero stratum is covered by \eqref{eq:clearing}.

Every degree from zero through twenty is a sum of two exponents in
\[
                        \{0,1,2,4,5,8,9,10\}.
\]
Successive quadratic rows compute $t_*^2,t_*^4,t_*^5,t_*^8,t_*^9,t_*^{10}$.
The coefficient table then expresses $F_4(t_*)$ as a quadratic polynomial
in this basis. At the base scale $a_*=1$, the same basis expresses
$\delta^{20}$, so no separate numerator power chain is needed.
At the other scale, the existing values $x=a_*\delta,x^2,x^4$
are extended by $x^5=xx^4$ and $x^{10}=(x^5)^2$, and the first equation
of \eqref{eq:clearing} uses $(x^{10})^2$.
The existing trace-clearing coordinate now represents $e\delta^2$.

The high-point approximants $G_2$ are affine expressions in the
retained $f^2$ and $w^2$ coordinates. They introduce no additional
field value. At $V_1$ the numerator is simply $z(V_1)$.
The shared new field powers are
\[
                       C_1^4,\ A_3^4,\ A_4^4,\ A_4^8,
\]
each with two rational coordinates. The fourth powers of
$B_0,B_2,B_3,B_4$ already occur in the Taylor circuits.
Each detector has five field factors and three private intermediate
products, using sixteen coordinates; the final product uses the
retained $W$ coordinates. The count is
\[
\begin{array}{lr}\toprule
\text{retained coordinates and contact/logarithm circuits}&239\\
\text{four shared field powers}&8\\
\text{five factors and intermediate products for two detectors}&32\\\midrule
\text{total coordinates}&279\\\midrule
\text{retained and contact/logarithm quadratic rows}&209\\
\text{shared power rows}&8\\
\text{factor and product rows}&36\\\midrule
\text{total quadratic rows}&253\\\bottomrule
\end{array}
\]

The producer clears denominators and primitive-normalizes each row
before squaring. The independent checker imports no producer. It
reconstructs the degree-twenty contact polynomial by forward differences
of twenty-one exact evaluations, checks all twenty roots and the full
quadratic basis, reconstructs every changed row and every affine Taylor
substitution, verifies all retained rows, and independently expands the
whole sum of squares by unordered pairs. It separately checks the
$\delta^2$ trace normalization, the global $G_2$ nonvanishing certificate,
and 70912 supplementary valuation cases. The infinite arithmetic
arguments are the proofs in the preceding sections.

The expanded quartic has 3693 terms before powers of $A_N$ belonging
to the same coordinate monomial are collected, and 3692 monomials
after collection. Its unique coefficient involving $A_N^2$ is
$A_N^2+1$, at the square of the actual scalar contact-product coordinate.
Every other coefficient is a nonzero constant or a nonzero constant
multiple of $A_N$. Consequently no support cancellation occurs for
positive $A_N$. The largest constant coefficient $C_0$ has 1520 bits,
the largest linear multiplier is $13118436242$, and exact comparison
gives
\[
 C_0>13118436242^2,\qquad
 80989^{90}+1<C_0<80989^{94}+1.
\]
Thus $H(P_N)=\max(C_0,A_N^2+1)$ for every $N\ge1$ and the asserted
exact height holds for $N\ge23$. This is a coefficient statement;
it does not make the defining threshold effective.

\begin{thebibliography}{99}
\bibitem{Dependencies} Project repository,
\emph{Arithmetic dependencies of the defining-polynomial constructions},
\nolinkurl{arithmetic_dependencies_20261108_v8.tex}.
\bibitem{Source} Project repository,
\emph{A family of quartics in 297 variables},
\nolinkurl{quartic297_20261109_v12.tex}, especially Sections 2--6.
\bibitem{Weights} Project repository,
\emph{Common local weights and quartics in 297 and 342 variables},
\nolinkurl{shared_local_weights_20261109_v12.tex}.
\bibitem{Even} Project repository,
\emph{An explicit quartic in 322 variables},
\nolinkurl{quartic322_20261109_v13.tex}.
\bibitem{Contacts} Project repository,
\emph{Height Bounds from Disjoint Sets of Rational Contacts},
\nolinkurl{contact_height_amplification_20261109_v13.tex}.
\bibitem{Aggregate} Project repository,
\emph{The height of an aggregate elliptic trace},
\nolinkurl{aggregate_trace_height_20261108_v9.tex}.
\bibitem{Selection} Project repository,
\emph{Prime selection with fixed congruences and separated reduction orders},
\nolinkurl{prime_selection_20261109_v12.tex}.
\end{thebibliography}
\end{document}
